\section{问题重述}
长江流域是我国淡水渔业最发达、生物多样性最丰富的区域之一，但长期受过度捕捞、工业废水排放、水域塑料污染以及葛洲坝、三峡大坝等水利工程的综合影响，鱼类洄游通道被阻断，渔业资源持续衰退，典型食物链结构发生退化。自2021年1月1日实施长江十年禁渔政策以来，长江流域生态环境已出现显著变化：局部水域水质改善，水草等水生植被逐步恢复；2025年监测数据显示，青、草、鲢、鳙四大家鱼资源量已恢复至禁渔前的约1.8倍，表明禁渔对基础资源和经济鱼类具有显著的修复作用。

但生态系统的恢复并非总是朝着良性方向发展。部分水域在禁渔后出现鱼类种群数量增长过快、过度摄食水生植物和浮游动物的现象，进而导致水体浑浊、底栖植被大面积减少，部分水域甚至形成大面积"水下荒漠"。长江江豚种群数量已回升至约1249头，2024年长江鲟自然繁殖产卵场重新被发现；而鳄雀鳝、巴西龟等外来入侵物种已在鄱阳湖等流域支流定殖并维持种群。十年禁渔政策确实取得了显著的资源恢复成效，但恢复后的生态系统未必更加健康，营养级失衡、江湖连通性恢复不足以及外来物种入侵扩散等问题依然突出。因此，有必要分别度量"资源恢复"与"系统健康"，并追踪两者在各营养级上的差异。\cite{report2022,report2023}。

题目要求依次回答以下五个子问题：
\begin{itemize}
\item 问题一：分析禁渔政策对水生植物、浮游动物及四大家鱼种群动态的影响。
\item 问题二：研究禁渔产生的生态收益能否向上传递至长江鲟和长江江豚等顶级捕食者。
\item 问题三：构建长江流域简化鱼类食物链模型并探讨其长期演化规律。
\item 问题四：判别三级食物链系统在不同承载力条件下的动力学行为（稳定、周期振荡或混沌）。
\item 问题五：将工业废水排放、塑料污染和外来物种入侵等胁迫因子纳入统一评价框架，综合评估十年禁渔对长江渔业生态系统的整体效应，并提出针对性的治理对策建议
\end{itemize}
各问统一采用以下时间口径：2021年为政策启动基准年，2022年为公报基准年，2025年为恢复效果对照年\cite{yangtzeA}。

